\subsection{HAPI FHIR Server Implementation}
\label{subsec:hapi_fhir_server}

The core component responsible for standardizing healthcare data within the
CFI-Care system is the FHIR Server. This server is implemented using the HAPI
FHIR JPA Server Starter framework \textbf{version 8.4.0}. It is built upon
the Java 17 Spring Boot ecosystem and provides a fully compliant implementation
of the FHIR Release 5 (R5) standard. This choice ensures that all clinical
data, regardless of its source, adheres to international interoperability
protocols and persists reliably within the database.

To tailor the server for the specific needs of this project, several important
configuration changes were applied to the default setup. To maintain a
lightweight and efficient runtime, advanced features that were not immediately
required were explicitly disabled. This includes the Clinical Reasoning module
and the CDS Hooks service. Similarly, bulk export and import operations were
deactivated to prevent unnecessary resource consumption. Conversely, the
OpenAPI documentation endpoints were enabled. This configuration allows
developers to inspect the API structure and test endpoints directly through a
web interface.

The persistence layer is powered by a PostgreSQL database connected via the
Hibernate JPA engine. Database connection pooling is managed by
\textbf{HikariCP 5.0.1}, which maintains a pool of up to 20 concurrent
connections, minimizing connection overhead under high clinical load. Specific
configurations were applied to the Hibernate settings to optimize performance.
The system uses the dedicated \texttt{HapiFhirPostgresDialect} to ensure
efficient SQL generation. Additionally, the database schema management is
configured to automatically update the table structures upon server startup,
ensuring that the database remains synchronized with the code.

% \begin{figure}[htbp]
%     \centering
%     \fbox{\begin{minipage}[c][10cm]{\textwidth}
%             \centering
%             \vspace{4cm}
%             \textbf{FIGURE INSTRUCTION:} \\
%             \textit{Draw the FHIR Server Configuration Block Diagram.} \\
%             \textbf{Content:} Create a main box labeled "HAPI FHIR Server (Java/Spring)". Inside it, show enabled modules: "FHIR R5 REST API", "JPA Persistence Layer", "OpenAPI/Swagger". Show disabled modules with a cross or greyed out: "Clinical Reasoning", "CDS Hooks", "Bulk Export". Show an arrow pointing down to "PostgreSQL Database". \\
%             \textbf{Recommended Tool:} Draw.io or Lucidchart.
%         \end{minipage}}
%     \caption{Configuration of the HAPI FHIR Server showing enabled and disabled modules.}
%     \label{fig:fhir_config}
% \end{figure}

The server is responsible for managing a wide array of clinical resources
necessary for the application. These include foundational entities such as
Patient and Practitioner, as well as clinical events like Encounters,
Procedures, and Observations. It also handles complex data types, including
MedicationRequests for prescriptions and Binary resources for storing documents
like PDF files.

For deployment, the server utilizes a secure multi-stage Docker build process.
The first stage uses Maven to compile the application code and generate the
necessary artifacts. The final runtime image is built upon a ``distroless''
Java 17 base image. This approach significantly reduces the final container
size and improves security by removing unnecessary operating system tools and
shells from the production environment. The repository also ships Helm charts
(\texttt{charts/hapi-fhir-jpaserver/}), making the server deployable to a
Kubernetes cluster without additional configuration, supporting horizontal
scaling for higher-load production environments.